\section{模型假设}

\begin{enumerate}
  \item 题目给出的负载、光伏和电价数据在其时间区间内代表平均功率，换算为电量后可用于区间能量平衡；
  \item 储能充电效率与放电效率均为 0.9，即一次完整往返效率为 0.81；容量与最大充放电能力在考察期内不衰减；
  \item 微网不能向外部电网售电，光伏或计划购电富余时允许弃电；
  \item 问题一日首、日末 SOC 均为 $6000\,\mathrm{kWh}$；问题二至四的 SOC 跨日连续，并以终端电量价值抑制有限时域末端放空；
  \item 计划购电量没有题目未给出的额外上限；只有实际供电不足部分计作紧急购电；
  \item 问题三每次调整相对上一版有效计划结算，下调部分保留 50\% 的原购电成本，上调部分按 1.5 倍价格计费；
  \item 问题四的未来实时价格并非事先完全可见，主方案用历史数据进行因果预测，实际费用按真实价格结算。
\end{enumerate}
